\section*{Apêndice}
\addcontentsline{toc}{section}{Apêndice}

\subsection*{Apêndice A: Levantamento de Requisitos}
\addcontentsline{toc}{subsection}{Apêndice A: Levantamento de Requisitos}
\hypertarget{ap:apendice}{}

Este apêndice apresenta o processo de levantamento de requisitos adotado para o desenvolvimento do SpinFlow. É importante explicitar que o contexto descrito neste material possui caráter didático e parcialmente simulado, sendo construído com o objetivo de apoiar o entendimento das técnicas de engenharia de requisitos aplicadas ao projeto.

Embora o domínio representado (academia de spinning, fluxo de check-in, gestão de turmas e recursos operacionais) seja plausível e coerente com cenários reais, não houve execução integral de um processo formal de elicitação em ambiente produtivo, com stakeholders reais e ciclos completos de validação empírica.

O levantamento de requisitos considerado para a elaboração desta especificação foi estruturado de forma didática e representativa, buscando refletir práticas compatíveis com projetos de software orientados à operação. Ainda que não tenha sido conduzido em um ambiente real com stakeholders ativos, o processo foi organizado com o propósito de simular a consolidação de informações relevantes sobre o contexto da academia, os problemas centrais do fluxo de \textit{check-in}, as necessidades dos usuários e os limites do recorte funcional adotado no \ProjetoNome.

Para compor essa base de trabalho, considera-se a utilização simulada e orientada por boas práticas de técnicas de engenharia de requisitos, com ênfase em:

\begin{itemize}[leftmargin=1.2cm]
\item \textbf{Entrevistas semiestruturadas}: utilizadas para levantar dores operacionais, expectativas de uso e percepções sobre os principais problemas ligados à reserva, à ocupação das aulas, à manutenção das bikes e à experiência do aluno.
\item \textbf{Sessões colaborativas do tipo \textit{JAD}}: utilizadas para consolidar entendimento entre as partes interessadas sobre o fluxo central do produto, suas prioridades e os limites do escopo.
\item \textbf{Análise documental}: utilizada para revisar artefatos do próprio projeto, regras de negócio, diagrama de classes, protótipos de interface e materiais de apoio relevantes para a especificação.
\item \textbf{Consolidação de artefatos}: utilizada para transformar as informações coletadas em um conjunto coerente de definições, escopo, regras e requisitos.
\end{itemize}

No contexto desta especificação, essas técnicas sustentam a identificação do \textit{check-in} como núcleo do produto, da Professora como frente de suporte operacional e do grupo de alunos como estrutura complementar de menor prioridade. Também sustentam a delimitação entre o que faz parte do recorte atual do sistema e o que deve permanecer como possibilidade de evolução futura. Em conjunto, esse processo permitiu consolidar uma especificação alinhada ao problema de negócio, ao domínio modelado e às decisões de escopo formalizadas neste documento.

\newpage

\subsection*{Apêndice B: Modelos de Análise}
\addcontentsline{toc}{subsection}{Apêndice B: Modelos de Análise}

\subsubsection*{B.1 Diagrama de Classes}
\begin{figure}[!h]
\centering
\includegraphics[width=\textwidth]{img/dc.png}
\caption{Diagrama de classes do projeto \ProjetoNome. Para preservar clareza visual, foram omitidos os fluxos e operações tradicionais de CRUD das entidades de cadastro.}
\end{figure}
\newpage

\subsubsection*{B.2 Diagramas de Casos de Uso}
\paragraph{Perfil Professora}
\begin{figure}[!h]
\centering
\includegraphics[width=0.9\textwidth]{img/dcup.png}
\caption{Diagrama de casos de uso do perfil Professora no projeto \ProjetoNome.}
\end{figure}

\newpage

\paragraph{Perfil Aluno}
\begin{figure}[!h]
\centering
\includegraphics[width=0.9\textwidth]{img/dcug_aluno.png}
\caption{Diagrama de casos de uso do perfil Aluno no projeto \ProjetoNome.}
\end{figure}
\newpage

\subsubsection*{B.3 Tabela de Casos de Uso}
\begin{longtable}{@{}>{\raggedright\arraybackslash}p{0.10\textwidth}>{\raggedright\arraybackslash}p{0.52\textwidth}>{\centering\arraybackslash}p{0.14\textwidth}>{\centering\arraybackslash}p{0.14\textwidth}@{}}
\toprule
\textbf{ID} & \textbf{Caso de Uso} & \textbf{Professora} & \textbf{Aluno} \\ \hline
\midrule
\endfirsthead
% \toprule
% \textbf{ID} & \textbf{Caso de Uso} & \textbf{Professora} & \textbf{Aluno} \
% \midrule
% \endhead



UC01 & Autenticar no Sistema & Sim & Sim \\ \hline
UC02 & Realizar Checkin & Sim & Sim \\ \hline
UC03 & Verificar Disponibilidade da Bike (\textit{include} de UC02) & --- & --- \\ \hline
UC04 & Visualizar Mapa da Sala (\textit{include} de UC02) & --- & --- \\ \hline
UC05 & Cancelar Checkin (\textit{extend} de UC02) & Sim, qualquer & Sim, próprio \\ \hline
UC06 & Gerenciar Alunos e Grupos & Sim & --- \\ \hline
UC07 & Gerenciar Bikes e Fabricantes & Sim & --- \\ \hline
UC08 & Gerenciar Manutenções & Sim & --- \\ \hline
UC09 & Cancelar Manutenção (\textit{extend} de UC08) & Sim & --- \\ \hline
UC10 & Gerenciar Salas e Turmas & Sim & --- \\ \hline
UC11 & Verificar Disponibilidade de Horário (\textit{include} de UC10) & --- & --- \\ \hline
UC12 & Associar Mix a Turma (\textit{extend} de UC10) & Sim & --- \\ \hline
UC13 & Gerenciar Mixes e Repertório & Sim & --- \\ \hline
UC14 & Consultar Agenda Semanal & Sim & Sim \\ \hline
UC15 & Consultar Mix e Repertório & Sim & Sim \\ \hline
UC16 & Consultar Histórico de Presença & --- & Sim \\ \hline
UC17 & Visualizar Dashboard & Sim & Sim \\ \hline
UC18 & Gerar Relatórios Gerenciais & Sim & --- \\ \hline
\bottomrule
\end{longtable}